\documentclass{beamer}
\usetheme{Berlin}

\title{Systems with a Small Number of
Variables}
\subtitle{Course: Complex Systems Modeling}
\author{Robert Flassig}
\institute{Brandenburg University of Applied Sciences}
\date{\today}

\begin{document}
\begin{frame}
\titlepage
\end{frame}

\begin{frame}
\frametitle{Outline}
\tableofcontents
\end{frame}
\section{Bifurcation Theory}

\begin{frame}
\frametitle{What is a Bifurcation?}

\begin{itemize}
    \item A \textbf{bifurcation} occurs when a small smooth change in a parameter
    causes a sudden qualitative change in the long-term behavior of a system.
    \item For a 1D autonomous system:
    \[
    \dot x = f(x,\mu)
    \]
    a bifurcation occurs when both:
    \begin{itemize}
        \item An equilibrium \(f(x^\ast,\mu)=0\) is created/destroyed or changes stability.
        \item The \textbf{linearization} \(f_x(x^\ast,\mu)\) becomes zero.
    \end{itemize}
    \item Bifurcations are the simplest building blocks of nonlinear phenomena
    (cf.\ Strogatz, \emph{Nonlinear Dynamics and Chaos}, 2nd ed., 2018, Ch.\,3).
\end{itemize}

\end{frame}

% ------------------------------------------
\begin{frame}
\frametitle{Why Bifurcations Cannot Occur in Linear Systems}

Linear system:
\[
\dot x = A x
\]

\begin{itemize}
    \item Eigenvalues of \(A\) determine stability.
    \item Changing parameters changes eigenvalues \emph{continuously}.
    \item But:
    \[
    \textbf{origin is always the only equilibrium.}
    \]
    \item No new fixed points can appear or disappear.
\end{itemize}

\textbf{Conclusion:}  
Bifurcations require \textbf{nonlinear} systems.

\end{frame}

% ============================================================
% Saddle-Node Bifurcation
% ============================================================

\subsection{Saddle-Node Bifurcation}

\begin{frame}
\frametitle{Example: Saddle--Node Bifurcation}

Canonical model:
\[
\dot x = \mu - x^2
\]

Equilibria:
\[
x^\ast = \pm \sqrt{\mu} \qquad (\mu>0)
\]

Stability:
\[
f_x = -2x^\ast \;\Rightarrow\;
\begin{cases}
x^\ast = +\sqrt{\mu}: \; f_x < 0 \Rightarrow \text{stable} \\
x^\ast = -\sqrt{\mu}: \; f_x > 0 \Rightarrow \text{unstable}
\end{cases}
\]

\textbf{At \(\mu=0\)}: Both equilibria collide and annihilate.

\end{frame}

\begin{frame}
\frametitle{Saddle--Node Phase Portraits}

\begin{center}
\fbox{\parbox{0.6\textwidth}{\centering \vspace{0.5cm} \textit{Saddle--node bifurcation diagram} \vspace{0.5cm}}}
\end{center}

\textbf{Qualitative change:}
\begin{itemize}
    \item \(\mu < 0\): No equilibria.
    \item \(\mu > 0\): Two equilibria (one stable, one unstable).
\end{itemize}

\end{frame}

% ============================================================
% Transcritical Bifurcation
% ============================================================

\subsection{Transcritical Bifurcation}

\begin{frame}
\frametitle{Example: Transcritical Bifurcation}

Canonical form:
\[
\dot x = \mu x - x^2.
\]

Equilibria:
\[
x_1^\ast = 0, \qquad x_2^\ast = \mu.
\]

Stability:
\[
f_x =
\begin{cases}
\mu & \text{at } x^\ast = 0 \\
-\mu & \text{at } x^\ast = \mu
\end{cases}
\]

\textbf{Interpretation:}  
The two fixed points exchange their stability at \(\mu = 0\).

\end{frame}

\begin{frame}
\frametitle{Transcritical Bifurcation Diagram}

\begin{center}
\fbox{\parbox{0.6\textwidth}{\centering \vspace{0.5cm} \textit{Transcritical bifurcation diagram} \vspace{0.5cm}}}
\end{center}

\textbf{Typical scenario:}
\begin{itemize}
    \item Two branches cross.
    \item Each branch is stable on one side and unstable on the other.
\end{itemize}

\end{frame}

% ============================================================
% Pitchfork
% ============================================================

\subsection{Pitchfork Bifurcation}

\begin{frame}
\frametitle{Supercritical Pitchfork Bifurcation}

Canonical form:
\[
\dot x = \mu x - x^3.
\]

Equilibria:
\[
x^\ast = 0,
\qquad
x^\ast = \pm\sqrt{\mu} \; (\mu > 0).
\]

Stability:
\[
f_x = \mu - 3x^2.
\]

\begin{itemize}
    \item At \(x=0\):
    \(\mu<0\) stable, \(\mu>0\) unstable.
    \item New stable equilibria appear when \(\mu>0\).
\end{itemize}

\end{frame}

\begin{frame}
\frametitle{Pitchfork Bifurcation Diagram}

\begin{center}
\fbox{\parbox{0.6\textwidth}{\centering \vspace{0.5cm} \textit{Pitchfork bifurcation diagram} \vspace{0.5cm}}}
\end{center}

\textbf{Symmetry breaking:}  
System transitions from one stable equilibrium to two stable equilibria.

\end{frame}

% ============================================================
% Linear Stability Criterion
% ============================================================

\subsection{Linear Stability at Bifurcation Points}

\begin{frame}
\frametitle{Failure of Linear Stability at Bifurcations}

At an equilibrium \(f(x^\ast,\mu)=0\), consider the linearization:
\[
\dot \eta = f_x(x^\ast,\mu)\eta.
\]

\textbf{Bifurcation condition:}
\[
f_x(x^\ast,\mu_c) = 0.
\]

Implication:
\begin{itemize}
    \item Linearization predicts \(\dot \eta = 0\) $\to$ \textbf{inconclusive}.
    \item Need higher-order terms $\to$ must use \textbf{nonlinear analysis}.
    \item Standard tool: \textbf{Taylor expand} near equilibrium:
    \[
    f(x^\ast+\eta,\mu) =
    f_x \eta + \frac{1}{2}f_{xx}\eta^2 + \dots
    \]
\end{itemize}

\end{frame}

% ============================================================
% Classification
% ============================================================

\begin{frame}
\frametitle{Classification of Codimension-1 Bifurcations}

Only three generic 1D continuous-time bifurcations:

\begin{enumerate}
    \item \textbf{Saddle--Node}: Two equilibria collide and annihilate.
    \item \textbf{Transcritical}: Two branches exchange stability.
    \item \textbf{Pitchfork}:
    \begin{itemize}
        \item Supercritical: one $\to$ three equilibria.
        \item Subcritical: opposite direction.
    \end{itemize}
\end{enumerate}

All others are non-generic or higher-dimensional.

\end{frame}

% ============================================================
% Exercises
% ============================================================

\subsection{Exercises}

\begin{frame}
\frametitle{Exercise: Identify Bifurcation Type}

Classify the bifurcation for each system:

\begin{enumerate}
    \item \( \dot x = \mu + x^2 \)
    \item \( \dot x = \mu x - x^2 \)
    \item \( \dot x = \mu x + x^3 \)
\end{enumerate}

\medskip
\textbf{Tasks:}
\begin{itemize}
    \item Find equilibria as function of \(\mu\).
    \item Determine stability.
    \item Identify the bifurcation type.
\end{itemize}

\end{frame}

\begin{frame}
\frametitle{Solutions}

\begin{enumerate}
    \item \( \dot x = \mu + x^2 \)\\
    $\to$ Saddle--Node (at \(\mu=0\)).

    \item \( \dot x = \mu x - x^2 \)\\
    $\to$ Transcritical (exchange of stability at \(\mu=0\)).

    \item \( \dot x = \mu x + x^3 \)\\
    $\to$ Subcritical Pitchfork (unstable branches).
\end{enumerate}

\textbf{Reference:} S.H. Strogatz, \emph{Nonlinear Dynamics and Chaos}, 2nd ed., Westview Press, 2018, Chapter 3.

\end{frame}



\end{document}